\scanpage{201}
\uline{Démonstration.} \deleted{Fixons une structure
\deleted{\unsure{conforme}} \inserted{\unsure{complexe}}
sur $X$, \unsure{invariante} par}

Démontrons d). \deleted{Soit $\Gamma_0 =$} Nous identifions
$\Gamma',\Gamma''$ à des ss-groupes de $\Gamma$, et posons $\Gamma_0=\Gamma'\cap\Gamma''$.
Soit \unsure{$\mathcal E^!$} le ss-groupe de $\mathcal E$ formé des $g\in\mathcal E$
tels que $\operatorname{int}(g).\Gamma_0$ soit $\pi$-conjugué de $\Gamma_0$
\unsure{i.e.} dont l'image dans $\Gamma$ \deleted{\illegible{1}} centralise $\Gamma_0$.
\deleted{\illegible{1}} On a $\mathcal E^!\supset\pi$, et $\mathcal E^!$ est donc l'image
inverse d'un sous-groupe $\Gamma^!$ de $\Gamma$, qui
centralise $\Gamma_0$ \deleted{\illegible{1}}, et qui contient $\Gamma'$ et $\Gamma''$
(car $\mathcal E^!$ contient le centralisateur de $\Gamma_0$ dans $\mathcal E$,
donc $\Gamma'$ et $\Gamma''$). Quitte à remplacer
$\Gamma$ par $\Gamma^!$, $\mathcal E$ par $\mathcal E^!$, OPS $\mathcal E=\mathcal E^!$, $\Gamma=\Gamma^!$ i.e.
que \unsure{$\Gamma_0\subset\mathrm{Centre}$|$\Gamma_0\in\mathrm{Centre}$} de $\Gamma$, et que \unsure{$\Gamma_0\hookrightarrow\mathcal E$}
\unsure{invariant} mod $\pi$-conjugaison par $\Gamma$.

\deleted{Après quoi choisissons une réalisation top.
$(X,S,a)$ de l'opération \deleted{ext.} de $\Gamma_0$ sur $\pi$,
où $a\in U=X-S$ -- donc $\pi\simeq\pi_1(U,a)$. D'après
le résultat (\unsure{b}) d'existence et d'unicité
des réalisations dans le cas cyclique, cela
est possible, et de plus on peut faire}

\scanpage{202}
\mspage{101}
\deleted{opérer $\Gamma'$ sur \unsure{$(X,S,a)$} (i.e. sur $(U,a)$) de façon
compatible avec son action sur $\pi_1(U,a)=\pi$, de
façon à prolonger l'action de \unsure{$\Gamma_0$} --
et de \unsure{$\overline{\mathrm m}$|$=$} pour $\Gamma''$. Mais il n'est pas clair
à priori
que les deux actions \deleted{\unsure{puissent}} soient compatibles
en le sens préc. \deleted{\unsure{du ss-groupe engendré}}
\inserted{que le \unsure{ss-groupe engendré}}
\deleted{\unsure{par} $\Gamma',\Gamma''$} $\subset\operatorname{Aut\,top}(U)$ soit aussi fini.}

\deleted{Considérons maintenant $X'=X/\Gamma_0$, $S'=S/\Gamma_0$,
et l'image $a'$ de $a$ dans $X'$ -- donc $a'\in U'=X'-S'$.
Alors \unsure{$\Gamma/\Gamma_0$|$\Gamma'/\Gamma_0$} opère \unsure{isotypiquement|isotopiquement} sur $(X',S',b)$}

\deleted{Utilisons $\pi^{\Gamma_0}=(1)$, on trouve que}

On va construire une section de $\mathcal E$
sur $\Gamma$ \unsure{tt} entier, ainsi. Soit $\widetilde\Gamma\subset\mathcal E$
le centralisateur de $\Gamma_0$ dans $\mathcal E$, \unsure{$\Rightarrow$}
\deleted{$\gamma\in\Gamma$, $g\in\mathcal E$ au-dessus de $\gamma$, considérons}
$\widetilde\Gamma\cap\pi=(1)$ (car cela signifie $\pi^{\Gamma_0}=(1)$)
\deleted{$\operatorname{int}(g).\Gamma_0$, $\exists!\ \alpha\in$}
\inserted{donc}
\deleted{je dis que} $\widetilde\Gamma\longrightarrow\Gamma$ injectif, je dis
qu'il est surjectif. En effet, soit
$\gamma\in\Gamma$, $g\in\mathcal E$ au dessus de $\gamma$, par hyp
\inserted{sur $\Gamma$}
$\exists\ \alpha\in\pi$ tel que $\operatorname{int}(g).\Gamma_0=\operatorname{int}(\alpha).\Gamma_0$
i.e. OPS $\operatorname{int}(g)\Gamma_0=\Gamma_0$, donc $g$ normalise
$\Gamma_0$, mais comme $\Gamma_0$ est central dans $\Gamma$
cela signifie que $g$ centralise $\Gamma_0$ i.e. $g\in\widetilde\Gamma$.

\scanpage{203}
Ainsi $\widetilde\Gamma$ \inserted{$\simeq\Gamma$} \unsure{est} un sous-groupe \uline{fini} de $\operatorname{Aut}(\pi)$,
\deleted{donc cyclique} contenant $\Gamma'$ et $\Gamma''$, cqfd.

\uline{NB} Si on n'avait pas au début supposé
$\Gamma$ fini, il serait vrai encore que $\Gamma',\Gamma''$
engendrent un ss-groupe $\widetilde\Gamma$ \inserted{$\subset\operatorname{Aut}(\pi)$} tel que
$\widetilde\Gamma\cap\pi=(1)$, mais cela nous fait une
belle jambe -- \unsure{on peut en trouver},
\deleted{p.ex. isom. à \unsure{$\operatorname{Autext}(\pi,i)$} tout entier,
(pour un \unsure{$i\in I(\pi)$} choisi: prendre le
\deleted{\illegible{2}}
centralisateur de \illegible{2} dans $\operatorname{Aut}(\pi)$)}

\scanpage{204}
\mspage{102}
19. Tour de Teichmüller.

Soit \inserted{$g\in\mathbb N$, et} $X_g$ une surface compacte \unsure{orientable} de
genre $g$, \deleted{\unsure{l'ensemble des orientations},}
\unsure{$[$|$]$} $(a_{g,i})_{i\in\mathbb N^*}=(a_{g,1},a_{g,2},\ldots)$ \unsure{$]$|$[$}
une suite de
pts distincts de $X_g$. On pose pour \unsure{$\forall\nu\in\mathbb N$}
\[
(1)\qquad S_{g,\nu}=\{\unsure{x_{g,i}|a_{g,i}}\mid 0\leq i\leq\nu-1\}
\qquad\text{donc }\nu=\operatorname{card}S_{g,\nu}
\]
\[
(2)\qquad U_{g,\nu}=X_g-S_{g,\nu}.
\]
\margin{NB on $a_{g,\nu}\in U_{g,\nu}$}
On a donc $S_{g,0}=\varnothing$, $U_{g,0}=X_g$; Les $S_{g,\nu}$ forment
une suite \inserted{str.} croissante de parties finies de $X_g$,
les $U_{g,\nu}$ une suite str. décroissante d'ouverts
de $X_g$. On prendra par la suite $U_{g,\nu}$ comme
la surface (orientée, ou orientable) type de
type $(g,\nu)$.

Soit
\[
(3)\qquad A_g=\operatorname{Aut}(X_g)
\]
le groupe des autom. de $X_g$, muni de la
top. de la conv. uniforme de $u$ et de son
inverse. On pose
\[
(4)\qquad A_{g,\nu}=\{u\in A_g\mid u(S_{g,\nu})=S_{g,\nu}\}\subset A_g
\]
ainsi $A_{g,\nu}$ est un sous-groupe de $A_g$. On
a un isomorphisme canonique $u\mapsto u|U_{g,\nu}$
\[
A_{g,\nu}\xrightarrow{\sim}\operatorname{Aut}(U_{g,\nu}).\qquad(5)
\]
\margin{\uline{NB} C'est \unsure{un isom.} \inserted{\illegible{3}}
topologiques, quand $A_{g,\nu}$ est
muni de la top. induite par $A_g$,
et $\operatorname{Aut}(U_{g,\nu})$ de la top. habituelle de la conv.
\unsure{comp.} \illegible{1} de $u$ et son inverse.}
On a \deleted{\unsure{un}} homomorphisme can. surjectif
\[
(6)\qquad A_{g,\nu}\longrightarrow\operatorname{Aut}(S_{g,\nu})
\bigl(\overset{\text{(can)}}{\simeq}\unsure{\mathfrak S_\nu}\bigr),
\]

\scanpage{205}
dont le noyau est noté par $A_{g,\nu}^!$
\[
(7)\qquad A_{g,\nu}^!=\{u\in A_g\mid u(a_{g,i})=a_{g,i},\ \forall i\in[0,\ldots,\nu-1]\}
\]
d'où suite exacte
\[
(8)\qquad 1\longrightarrow A_{g,\nu}^!\longrightarrow A_{g,\nu}
\longrightarrow\mathfrak S_\nu\longrightarrow1.
\]
Soit $A_{g,\nu}^{\circ}$ la composante neutre du groupe
$A_{g,\nu}$, on a donc
\[
(9)\qquad A_{g,\nu}^{\circ}\bigl(=(A_{g,\nu}^!)^{\circ}\bigr)\subset A_{g,\nu}^!.
\]
Posons
\[
(10)\qquad \Gamma_{g,\nu}=A_{g,\nu}/A_{g,\nu}^{\circ}=\pi_0(A_{g,\nu})
\quad\text{(groupe de Teichmüller de type }g,\nu\text{)}
\]
\margin{On pose aussi $\Gamma_g=\Gamma_{g,0}(=\Gamma_{g,0}^!)$}
\[
(11)\qquad \Gamma_{g,\nu}^!=A_{g,\nu}^!/A_{g,\nu}^{\circ}=\pi_0(A_{g,\nu}^!).
\]
La suite exacte (8) donne donc une suite
exacte
\[
(12)\qquad 1\longrightarrow\Gamma_{g,\nu}^!\longrightarrow\Gamma_{g,\nu}
\longrightarrow\mathfrak S_\nu\longrightarrow1.
\]
On a des homéomorphismes canoniques
\[
(13)\qquad A_g/A_{g,\nu}\simeq
\begin{gathered}
\inserted{\operatorname{Sym}^{\nu}(X)^*}\quad
\text{ouvert \inserted{du} \deleted{\unsure{$X_g$}} produit symé-}\\
\text{trique }\inserted{\unsure{\nu^{\text{ième}}}}\operatorname{Sym}^{\nu}(X_g)\text{ de }X_g\\
\text{formé des parties finies}\\
\text{de cardinal }\nu\\
=\unsure{\mathcal P_\nu(X_g)}
\end{gathered}
\]
\margin{NB les pts \unsure{des espaces} \illegible{2}
des premiers membres de (13), (14)
\unsure{s'en vont vers} les systèmes
de pts $\{\unsure{a_{g,0}},a_{g,1},\ldots,a_{g,\nu-1}\}$ ou
$(a_{g,0},\ldots,a_{g,\nu-1})$.}
\[
(14)\qquad A_g/A_{g,\nu}^!\simeq
\begin{gathered}
\inserted{(X_g^{\nu})^*}\quad\text{ouvert de }X_g^{\nu}\text{ formé des}\\
\nu\text{-uples de pts distincts}\\
\simeq\operatorname{Mon}(\deleted{\unsure{S}_g}I_\nu,X)\quad\text{(où }I_\nu\\
=\{0,\ldots,\nu-1\}\text{).}
\end{gathered}
\]

\scanpage{206}
\mspage{103}
Cet homéomorphisme \unsure{respectant} les actions
naturelles de $\mathfrak S_\nu\simeq A_{g,\nu}/A_{g,\nu}^!$.

Les $A_{g,\nu}^!$ pour $\nu$ variable forment une
suite décroissante de groupes de $A_g$
\[
(15)\qquad A_g\unsure{=|\supset}A_{g,0}^!\supset A_{g,1}^!\supset A_{g,2}^!
\supset\cdots\supset A_{g,\nu}^!\supset\cdots
\]
et les homomorphismes correspondants
entre espaces homogènes \inserted{de $A_g$} s'insèrent dans le
diagramme commutatif \unsure{d'homéo} \illegible{1}

(16)\quad
\diagram{\begin{tikzcd}[ampersand replacement=\&]
 A_g/A_{g,\nu}^! \arrow[r,"\sim"] \arrow[d] \& \operatorname{Mon}(I_\nu,X) \arrow[d] \\
 A_g/A_{g,\nu'}^! \arrow[r,"\sim"] \& \operatorname{Mon}(I_{\nu'},X)
\end{tikzcd}}
\quad $0\leq\nu'\leq\nu$
\margin{\unsure{fonct.}}
où la deuxième flèche verticale est celle de
restriction de $I_\nu$ à $I_{\nu'}\subset I_\nu$.

Pour $\nu=1$ \inserted{on a $A_{g,1}^!=A_{g,1}$, et} l'isomorphisme (13) (ou au choix,
(14)) s'écrit
\[
(17)\qquad A_g/A_{g,1}\simeq X_g
\]
(homéomorphisme compatible avec \inserted{les} actions de $A_g$)

D'ailleurs si à $\forall x\in X_g$ on associe son
stabilisateur $A_{g,x}$ dans $A_g$, on trouve une
application évidemment surjective

\scanpage{207}
\[
(18)\qquad X_g\longrightarrow
\begin{gathered}
\text{ensemble des conjugués du sous-}\\
\text{groupe }A_{g,1}\text{ de }A_g\\
\simeq A_g/\operatorname{Norm}_{A_g}(A_{g,1})
\end{gathered}
\]
qui s'identifie via (17) à l'application canonique
sur les espaces homogènes
\[
 A_g/A_{g,1}\longrightarrow A_g/\operatorname{Norm}_{A_g}(A_{g,1})
\]
déduite de l'inclusion $A_{g,1}\subset\operatorname{Norm}_{A_g}(A_{g,1})$,
et on voit de suite que (18) est
bijective i.e. que \deleted{\illegible{1}}
\[
(19)\qquad A_{g,1}=\operatorname{Norm}_{A_g}(A_{g,1}),
\]
plus généralement pour $\forall\nu$, on a
\[
(20)\qquad A_{g,\nu}=\operatorname{Norm}_{A_g}(A_{g,\nu})
=\operatorname{Norm}_{A_g}(A_{g,\nu}^!)
\]
ce qui signifie que les applications canoniques
de $A_g$-ensembles homogènes
\deleted{$x=\{x_1,\ldots,x_\nu\}\longmapsto A_{g,x}=\bigcap$}
\deleted{(23)}
\[
\mathcal P_\nu(X_g)\longrightarrow\text{ens. des conjugués de }A_{g,\nu}
\]
\[
(21)\qquad S\longmapsto\text{stabilisateur }A_{g,S}\text{ de }S
\]
et
\[
(22)\qquad\unsure{\vert}\quad S\longmapsto A_{g,S}^!
\]
sont non seulement surjectives, mais $\overline{\mathrm m}$ bijectives.
Cela provient du fait que l'on retrouve
$S$ en termes de $A_{g,S}^!$ \inserted{\unsure{ou de $A_{g,S}$}}
\deleted{\unsure{et a fortiori en termes} \unsure{du groupe} $A_{g,S}$, comme}

\scanpage{208}
\mspage{104}
\[
(23)\qquad S=\{x\in X_g\mid ux=x\quad\forall u\in A_{g,S}^!\}
=X_g^{(A_{g,S}^!)}
\]
\[
(24)\qquad\begin{aligned}
S&=\{x\in X_g\mid A_{g,S}.x\text{ fini}\}\\
 &=\{x\in X_g\mid A_{g,S}.x\ne X_g\}
\end{aligned}
\]
\margin{$\forall S\in\mathcal P_\nu(X_g)$}

On a aussi des applications canoniques
d'espaces homogènes sous $A_g$
\[
(25)\qquad S\longmapsto A_{g,S}^{\circ}
\]
\[
\mathcal P_\nu(X_g)\longrightarrow
\begin{gathered}
\text{ensemble des conjugués}\\
\text{des }A_{g,\nu}^{\circ}\text{ dans }A_g\\
\simeq A_g/\operatorname{Norm}(A_{g,\nu}^{\circ})
\end{gathered}
\]
qui est non seulement surjective, mais encore
injective, car on a la relation suivante, qui
renforce \inserted{(23) et} (24), $\forall S\in\mathcal P_\nu(X_g)$
\[
(26)\qquad\begin{aligned}
S&=\{x\in X_g\mid A_{g,S}^{\circ}.x=\{x\}\}=X_g^{A_{g,S}^{\circ}}\\
 &=\{x\in X_g\mid A_{g,S}^{\circ}.x\text{ fini}\}\\
 &=\{x\in X_g\mid A_{g,S}^{\circ}.x\ne X_g\}.
\end{aligned}
\]
Ainsi on a \deleted{\illegible{1}} pour $\forall\nu$
\[
(27)\qquad A_{g,\nu}=\operatorname{Norm}_{A_g}(\unsure{A_{g,\nu}|A_{g,x}})
=\operatorname{Norm}_{A_g}(A_{g,\nu}^!)=\operatorname{Norm}_{A_g}(A_{g,\nu}^{\circ}).
\]

Soit $G$ un groupe topologique, muni
d'une classe de conjugaison \unsure{$\mathcal E$|$\mathcal X$} de ss-groupes,
\deleted{Soit} $G_1$ dans cette classe. On dit que

\scanpage{209}
$G,\unsure{\mathcal E|\mathcal X}$ est un \deleted{\unsure{donnée de}} \inserted{couple de} Teichmüller
de type $g$, s'il existe un iso. de groupes
topologiques $G\simeq A_g$, transformant
\unsure{$\mathcal E$|$\mathcal X$} en la classe de conjugaison de $A_1$.
\inserted{\unsure{Il revient au $\overline{\mathrm m}$ de dire que}}
\deleted{On voit alors que $G_1=\operatorname{Norm}_G(G_1)$}
\deleted{donc que} \unsure{$\mathcal E$|$\mathcal X$}, \deleted{\unsure{comme}}
\inserted{muni de la top. d'}espace homogène
sous $G$, est \deleted{homéom. à $G/G_1$} une
surface compacte orientable de genre $g$,
\deleted{\illegible{2}} soit \unsure{$\mathcal X$|$\mathfrak X$}, et que
\deleted{\illegible{3}} \inserted{l'application}
naturelle
\[
(28)\qquad G\longrightarrow\operatorname{Aut}(\mathcal X)
\]
est un homéomorphisme de groupes topologiques.
On voit alors que \inserted{\unsure{par} $(G,\unsure{\mathcal E|\mathcal X})\longrightarrow\mathcal X$ de}
la catégorie des
couples de Teichmüller de type $g$, vers
celle des surfaces compactes \unsure{orientées|orientables}, est
une équivalence de catégories. Il en résulte
que \deleted{\unsure{tout auto}} pour un autom. \inserted{\unsure{$u$}} du
groupe top. $G$, $u$ est intérieur ssi il

\scanpage{210}
\mspage{105}
conserve la classe \unsure{$\mathcal E$|$\mathcal X$} i.e. ssi $u(G_1)$ est conjugué
de $G_1$. D'ailleurs le centre de $G$ est trivial, \unsure{i.e.}

On peut donner une description analogue
pour la catégorie des surfaces \unsure{orientées|orientables} de
genre $g$ à $\nu$ trous (i.e. homéomorphes
à $U_{g,\nu}$) en termes d'un groupe topo-
logique $G_\nu(\simeq A_{g,\nu})$ et d'une \inserted{classe de conj. \unsure{$U_\nu$|$\widetilde U_\nu$} de}
ss-groupes $H_\nu$ de celui-ci $(\simeq A_{g,\nu}\cap A_{g,\nu+1})$, avec les
conditions que $(G_\nu,\{H_\nu\})$ soit isomorphe
à $(A_{g,\nu},A_{g,\nu}\cap A_{g,\nu+1}\unsure{)|\}}$, ou encore
que $U_\nu$ \inserted{(espace top. \unsure{d'esp. homog.})} soit une surface \unsure{orientée|orientable} de
type $(g,\nu)$, et que l'application can.
\[
(29)\qquad G\longrightarrow\operatorname{Aut}(U_\nu)
\]
soit un homéomorphisme de groupes
topologiques. On trouve alors une
équivalence entre la catégorie des
couples de Teichmüller $(G_\nu,\{H_\nu\})$
de type $(g,\nu)$, et celle des surfaces
orientables de type $(g,\nu)$. On trouve encore
